\chapter{Linear Algebra}
\label{ch:linalg}

\begin{tcolorbox}[feynbox={Sky}{BBg},title={Before and After}]
\textbf{Before this chapter you need:} functions and composition
(\S\ref{sec:functions}, \S\ref{sec:composition}), linear functions and slope
(\S\ref{sec:linear}), gradients and Jacobians
(\S\ref{sec:gradient}, \S\ref{sec:jacobian}).\\
\textbf{After it you will understand:} what a vector and a matrix are, why matrix
multiplication is defined the way it is, what a norm measures, and what
$\norm{\nabla D}_2 = 1$ demands of the critic --- the last piece needed before the
gradient penalty of Chapter~\ref{ch:gans} can be read as a sentence rather than a
formula.
\end{tcolorbox}

Chapter~\ref{ch:calculus2} repeatedly collected numbers into lists and grids and
promised a proper account later. This is that account. Neural networks are, at
bottom, alternating stacks of matrix multiplications and simple non-linear
functions, so the vocabulary here is the vocabulary in which every model in this
thesis is written.

%% ─────────────────────────────────────────────────────────────────────────────
\section{Scalars, Vectors and Matrices}
\label{sec:vectors}

\situation{
A shop tracks three products. Monday's sales are 4, 7 and 2 units. You could
write three separate numbers, or you could write one row, $[4, 7, 2]$, and treat
Monday as a single object. Add Tuesday, Wednesday and Thursday and you have a
table --- one row per day, one column per product.
}

\problem{
Handling each number separately means writing a separate line of algebra for each,
and a network layer would need a thousand such lines. Treating the collection as
one object lets a single symbol stand for the lot and a single operation act on
all of it at once.
}

\workedarith{
\textbf{Vectors.} Monday $\mathbf{u} = [4, 7, 2]$, Tuesday $\mathbf{v} = [1, 0, 5]$.

Addition is entry by entry:
\[
  \mathbf{u} + \mathbf{v} = [4+1,\; 7+0,\; 2+5] = [5, 7, 7] .
\]
Total units sold across the two days, per product. The interpretation survives the
operation, which is the sign that the operation is the right one.

Scaling multiplies every entry:
\[
  3\mathbf{u} = [12, 21, 6] ,
\]
what three identical Mondays would produce.

\medskip
\textbf{The dot product.} Suppose prices are $\mathbf{p} = [10, 2, 5]$ per unit.
Monday's revenue is
\[
  \mathbf{u} \cdot \mathbf{p} = 4 \times 10 + 7 \times 2 + 2 \times 5
  = 40 + 14 + 10 = 64 .
\]
Multiply matching entries, add the results. One number out of two lists.

\medskip
\textbf{Matrices.} Stack the two days:
\[
  A = \begin{bmatrix} 4 & 7 & 2 \\ 1 & 0 & 5 \end{bmatrix} ,
\]
a $2 \times 3$ matrix --- 2 rows, 3 columns, always in that order. The entry
$A_{23} = 5$ sits in row 2, column 3: Tuesday's sales of product 3.
}

\formaldef{
A \textbf{scalar} is a single number. A \textbf{vector} is an ordered list of
numbers; the number of entries is its \textbf{dimension}. A \textbf{matrix} is a
rectangular array, of \textbf{size} $m \times n$ for $m$ rows and $n$ columns.

For vectors of equal dimension, addition and scalar multiplication act entrywise.
The \textbf{dot product} of $\mathbf{u}, \mathbf{v} \in \R^n$ is
\[
  \mathbf{u} \cdot \mathbf{v} = \sum_{i=1}^{n} u_i v_i ,
\]
a scalar. Entry $A_{ij}$ of a matrix is in row $i$, column $j$.
}

\analogybreak{
The spreadsheet picture makes a matrix look like passive storage. Most matrices in
this thesis are not data at all --- they are \emph{actions}. A layer's weight
matrix is a thing that transforms its input, and the next section defines
multiplication precisely so that this reading is the correct one.
}

\whyhere{
A window of 128 daily returns is a vector in $\R^{128}$; a batch of 64 such
windows is a $64 \times 128$ matrix. The noise fed to the QuantGAN generator is a
vector in $\R^{100}$. These are the literal objects the code moves around.
}

%% ─────────────────────────────────────────────────────────────────────────────
\section{Matrix Multiplication as Composition}
\label{sec:matmul}

\situation{
Two currency conversions in sequence: pounds to euros, then euros to dollars. Each
is a multiplication. Doing both is a single conversion from pounds to dollars, and
its rate is the product of the two rates. The composite is itself a conversion of
the same kind.
}

\problem{
Layers of a network apply in sequence, exactly like the conversions. If each layer
is a matrix, the composition of two layers should again be a matrix --- otherwise
composing them would produce some new kind of object and the theory would not
close. Matrix multiplication is \emph{defined} so that it does close.
}

\workedarith{
\textbf{Matrix times vector first.} Let
\[
  A = \begin{bmatrix} 2 & 0 \\ 1 & 3 \end{bmatrix} ,
  \qquad \mathbf{x} = \begin{bmatrix} 4 \\ 5 \end{bmatrix} .
\]
Each output entry is the dot product of a row of $A$ with $\mathbf{x}$:
\[
  A\mathbf{x} = \begin{bmatrix} 2 \times 4 + 0 \times 5 \\ 1 \times 4 + 3 \times 5 \end{bmatrix}
  = \begin{bmatrix} 8 \\ 19 \end{bmatrix} .
\]

\medskip
\textbf{Now compose.} Let
\[
  B = \begin{bmatrix} 1 & 1 \\ 0 & 2 \end{bmatrix} .
\]
Apply $A$ then $B$ to $\mathbf{x}$:
\[
  B(A\mathbf{x}) = B\begin{bmatrix} 8 \\ 19 \end{bmatrix}
  = \begin{bmatrix} 1 \times 8 + 1 \times 19 \\ 0 \times 8 + 2 \times 19 \end{bmatrix}
  = \begin{bmatrix} 27 \\ 38 \end{bmatrix} .
\]

\medskip
\textbf{The composite matrix.} $BA$ has entry $(i,j)$ equal to the dot product of
row $i$ of $B$ with column $j$ of $A$:
\[
  BA = \begin{bmatrix} 1 & 1 \\ 0 & 2 \end{bmatrix}
       \begin{bmatrix} 2 & 0 \\ 1 & 3 \end{bmatrix}
     = \begin{bmatrix} 1\cdot2 + 1\cdot1 & 1\cdot0 + 1\cdot3 \\
                       0\cdot2 + 2\cdot1 & 0\cdot0 + 2\cdot3 \end{bmatrix}
     = \begin{bmatrix} 3 & 3 \\ 2 & 6 \end{bmatrix} .
\]
Test it on $\mathbf{x}$:
\[
  (BA)\mathbf{x} = \begin{bmatrix} 3 \times 4 + 3 \times 5 \\ 2 \times 4 + 6 \times 5 \end{bmatrix}
  = \begin{bmatrix} 27 \\ 38 \end{bmatrix} . \quad\checkmark
\]
Same answer. One matrix now does the work of two.

\medskip
\textbf{Order matters.} Reversing gives
\[
  AB = \begin{bmatrix} 2 & 0 \\ 1 & 3 \end{bmatrix}
       \begin{bmatrix} 1 & 1 \\ 0 & 2 \end{bmatrix}
     = \begin{bmatrix} 2 & 2 \\ 1 & 7 \end{bmatrix} \neq BA .
\]
Just as painting-then-boxing differed from boxing-then-painting in
\S\ref{sec:composition}.
}

\formaldef{
For $A$ of size $m \times k$ and $B$ of size $k \times n$, the product $AB$ has
size $m \times n$ with
\[
  (AB)_{ij} = \sum_{l=1}^{k} A_{il} B_{lj} .
\]
The inner dimensions must match: an $m \times k$ can only multiply a $k \times n$.

Matrix multiplication is \textbf{associative}, $(AB)C = A(BC)$, and
\textbf{distributive} over addition, but \textbf{not commutative}: $AB \neq BA$
in general.
}

\analogybreak{
Composition of matrices captures composition of \emph{linear} maps only. A
network built purely from matrices would collapse: stack fifty of them and the
product is a single matrix, so fifty layers would have exactly the expressive
power of one. The non-linear activation between layers is what prevents the
collapse, and it is the entire reason deep networks are deeper than shallow ones.
}

\whyhere{
A dense layer computes $W\mathbf{x} + \mathbf{b}$ --- a matrix multiplication and a
shift. The GRU cells in TimeGAN and the convolutions in QuantGAN and FinGAN are
structured matrix multiplications. The parameter counts reported for this
project's models ($9{,}744$ for the TimeGAN embedder, $11{,}400$ for its generator)
are counts of entries in these matrices.
}

%% ─────────────────────────────────────────────────────────────────────────────
\section{Transpose, Identity and Inverse}
\label{sec:inverse}

\situation{
Undoing things. Multiplying by 5 is undone by dividing by 5, and doing both leaves
the number untouched. Is there a matrix that undoes another, and a matrix that
does nothing at all?
}

\problem{
Backpropagation runs a network backwards, and solving a linear system means
undoing a transformation. Both need an account of reversal --- including an
account of when reversal is impossible, which turns out to be the more informative
case.
}

\workedarith{
\textbf{Transpose:} flip rows and columns.
\[
  A = \begin{bmatrix} 4 & 7 & 2 \\ 1 & 0 & 5 \end{bmatrix} ,
  \qquad
  A^{\!\top} = \begin{bmatrix} 4 & 1 \\ 7 & 0 \\ 2 & 5 \end{bmatrix} .
\]
A $2 \times 3$ becomes $3 \times 2$; entry $(i,j)$ moves to $(j,i)$.

\medskip
\textbf{Identity:} ones on the diagonal, zeros elsewhere.
\[
  I = \begin{bmatrix} 1 & 0 \\ 0 & 1 \end{bmatrix} ,
  \qquad
  I \begin{bmatrix} 4 \\ 5 \end{bmatrix}
  = \begin{bmatrix} 1 \cdot 4 + 0 \cdot 5 \\ 0 \cdot 4 + 1 \cdot 5 \end{bmatrix}
  = \begin{bmatrix} 4 \\ 5 \end{bmatrix} . \quad\checkmark
\]

\medskip
\textbf{Inverse.} For $A = \begin{bmatrix} 2 & 0 \\ 1 & 3 \end{bmatrix}$, the
inverse is
\[
  A^{-1} = \begin{bmatrix} 0.5 & 0 \\ -1/6 & 1/3 \end{bmatrix} .
\]
Check:
\[
  A^{-1}A = \begin{bmatrix} 0.5 \cdot 2 + 0 \cdot 1 & 0.5 \cdot 0 + 0 \cdot 3 \\
                            -\tfrac16 \cdot 2 + \tfrac13 \cdot 1 & -\tfrac16 \cdot 0 + \tfrac13 \cdot 3 \end{bmatrix}
  = \begin{bmatrix} 1 & 0 \\ -\tfrac13 + \tfrac13 & 1 \end{bmatrix}
  = \begin{bmatrix} 1 & 0 \\ 0 & 1 \end{bmatrix} = I . \quad\checkmark
\]

\medskip
\textbf{When it fails.} Take
$C = \begin{bmatrix} 1 & 2 \\ 2 & 4 \end{bmatrix}$, whose second row is twice the
first. Then
\[
  C\begin{bmatrix} 2 \\ -1 \end{bmatrix} = \begin{bmatrix} 0 \\ 0 \end{bmatrix}
  \quad\text{and}\quad
  C\begin{bmatrix} 0 \\ 0 \end{bmatrix} = \begin{bmatrix} 0 \\ 0 \end{bmatrix} .
\]
Two different inputs, identical output. No rule could recover which one you
started from, so no inverse exists. Information was destroyed --- the same failure
mode as multiplying an equation by zero in \S\ref{sec:variables}.
}

\formaldef{
The \textbf{transpose} $A^{\!\top}$ satisfies $(A^{\!\top})_{ij} = A_{ji}$, and
$(AB)^{\!\top} = B^{\!\top}A^{\!\top}$ --- the order reverses.

The \textbf{identity} $I_n$ has $I_{ii}=1$ and $I_{ij}=0$ for $i \neq j$; it
satisfies $AI = IA = A$.

A square $A$ is \textbf{invertible} if some $A^{-1}$ has $AA^{-1} = A^{-1}A = I$.
Such an $A^{-1}$ is unique when it exists. A matrix with no inverse is
\textbf{singular}, and $A$ is singular exactly when some non-zero vector satisfies
$A\mathbf{x} = \mathbf{0}$.
}

\analogybreak{
Invertibility is a yes/no question in theory and a matter of degree in practice.
A matrix can be technically invertible while being so close to singular that
computing $A^{-1}$ amplifies rounding error catastrophically. Numerical work
therefore speaks of \emph{conditioning} --- how nearly singular a matrix is ---
rather than treating invertibility as a clean binary.
}

\whyhere{
The transpose appears throughout backpropagation: if a layer's forward pass
multiplies by $W$, the corresponding backward pass multiplies by $W^{\!\top}$.
Singularity is also the right way to understand mode collapse --- a generator that
maps many different noise inputs to the same output has destroyed information in
exactly the way $C$ above does, and no amount of downstream processing can recover
the diversity that was lost.
}

%% ─────────────────────────────────────────────────────────────────────────────
\section{The Determinant}
\label{sec:determinant}

\situation{
Take the unit square --- corners at $(0,0)$, $(1,0)$, $(0,1)$, $(1,1)$, area 1 ---
and apply a matrix to each corner. You get a parallelogram. Its area tells you the
factor by which this transformation inflates or shrinks regions.
}

\problem{
We need a single number answering ``does this transformation destroy
information?''. Squashing a two-dimensional region flat onto a line reduces its
area to zero, and that is exactly the situation in which the map cannot be
undone.
}

\workedarith{
\textbf{A stretch.} $A = \begin{bmatrix} 2 & 0 \\ 0 & 3 \end{bmatrix}$ sends
$(1,0) \mapsto (2,0)$ and $(0,1) \mapsto (0,3)$. The unit square becomes a
$2 \times 3$ rectangle of area 6.

Formula for $2 \times 2$: $\det\begin{bmatrix} a & b \\ c & d\end{bmatrix} = ad - bc$.
Here $2 \times 3 - 0 \times 0 = 6$. \checkmark

\medskip
\textbf{A shear.} $B = \begin{bmatrix} 1 & 1 \\ 0 & 1 \end{bmatrix}$ sends
$(1,0) \mapsto (1,0)$ and $(0,1) \mapsto (1,1)$. The square becomes a
parallelogram with base 1 and height 1 --- area still 1. And
$\det B = 1 \times 1 - 1 \times 0 = 1$. \checkmark A shear slides without
changing area.

\medskip
\textbf{A collapse.} $C = \begin{bmatrix} 1 & 2 \\ 2 & 4 \end{bmatrix}$ from
\S\ref{sec:inverse} sends $(1,0) \mapsto (1,2)$ and $(0,1) \mapsto (2,4)$. Both
images lie on the same line through the origin --- $(2,4)$ is exactly twice
$(1,2)$. The square is flattened onto a line, area zero. And
$\det C = 1 \times 4 - 2 \times 2 = 0$. \checkmark

The determinant being zero and the inverse failing to exist are the same fact.
}

\formaldef{
The \textbf{determinant} $\det A$ of a square matrix is the signed factor by
which it scales volume. For $2 \times 2$,
\[
  \det\begin{bmatrix} a & b \\ c & d \end{bmatrix} = ad - bc .
\]
Key properties:
\begin{itemize}
  \item $\det A = 0$ if and only if $A$ is singular;
  \item $\det(AB) = \det A \cdot \det B$;
  \item a negative determinant means orientation is flipped, like a mirror.
\end{itemize}
}

\analogybreak{
The volume picture is exact but the computation does not scale. Evaluating a
determinant by the schoolbook cofactor expansion costs on the order of $n!$
operations --- for $n = 20$ that is $10^{18}$, beyond reach. Practical algorithms
use factorisations costing about $n^3$ instead. Never mistake a clean definition
for a usable algorithm.
}

\whyhere{
Determinants govern how probability densities transform under an invertible map,
via the Jacobian determinant --- the mechanism underpinning normalising flows,
which Chapter~12 (Generative Models) contrasts with GANs. GANs deliberately avoid
this machinery, which is what frees them from any invertibility requirement and
also why they provide no likelihood.
}

%% ─────────────────────────────────────────────────────────────────────────────
\section{Norms}
\label{sec:norms}

\situation{
You are at a street corner and a friend is three blocks east and four blocks
north. How far away is the friend? If you can fly, five blocks --- straight line.
If you must follow streets, seven --- three plus four. Both are honest answers to
``how far''; they differ because they measure different things.
}

\problem{
``How big is this vector?'' has no single answer, and the gradient penalty in
WGAN-GP asks precisely that question about $\nabla D$. Before that constraint can
be understood, the notion of size must be pinned down --- including which one is
meant.
}

\workedarith{
Take $\mathbf{v} = [3, 4]$.

\medskip
\textbf{Euclidean ($2$-norm) --- straight line:}
\[
  \norm{\mathbf{v}}_2 = \sqrt{3^2 + 4^2} = \sqrt{9 + 16} = \sqrt{25} = 5 .
\]

\textbf{Manhattan ($1$-norm) --- along the streets:}
\[
  \norm{\mathbf{v}}_1 = |3| + |4| = 7 .
\]

\textbf{Maximum ($\infty$-norm) --- the largest single component:}
\[
  \norm{\mathbf{v}}_\infty = \max(|3|, |4|) = 4 .
\]

\medskip
Three sizes for one vector: $5$, $7$, $4$. Always state which.

\medskip
\textbf{A three-entry example.} $\mathbf{w} = [1, -2, 2]$:
\[
  \norm{\mathbf{w}}_2 = \sqrt{1 + 4 + 4} = \sqrt{9} = 3, \qquad
  \norm{\mathbf{w}}_1 = 1 + 2 + 2 = 5, \qquad
  \norm{\mathbf{w}}_\infty = 2 .
\]

\medskip
\textbf{Scaling.} Doubling a vector doubles every norm:
$\norm{2\mathbf{w}}_2 = \sqrt{4 + 16 + 16} = \sqrt{36} = 6 = 2 \times 3$. \checkmark

\medskip
\textbf{Unit length.} Dividing by the $2$-norm gives length exactly 1:
$\mathbf{v}/5 = [0.6, 0.8]$, and $\sqrt{0.36 + 0.64} = 1$. \checkmark
This is the vector that won the steepest-direction contest in
\S\ref{sec:gradient} --- now recognisable as the gradient rescaled to unit length.
}

\formaldef{
A \textbf{norm} $\norm{\cdot}$ assigns a non-negative size to each vector, subject
to three requirements:
\begin{enumerate}
  \item $\norm{\mathbf{v}} = 0$ only when $\mathbf{v} = \mathbf{0}$;
  \item $\norm{c\mathbf{v}} = |c|\,\norm{\mathbf{v}}$ for any scalar $c$;
  \item $\norm{\mathbf{u} + \mathbf{v}} \leq \norm{\mathbf{u}} + \norm{\mathbf{v}}$
        (the triangle inequality: no detour is shorter than going direct).
\end{enumerate}
The \textbf{$p$-norm} is
\[
  \norm{\mathbf{v}}_p = \left(\sum_i |v_i|^p\right)^{1/p} ,
\]
giving the Manhattan norm at $p=1$, the Euclidean at $p=2$, and the maximum norm
in the limit $p \to \infty$. Unsubscripted $\norm{\cdot}$ conventionally means the
Euclidean norm.
}

\analogybreak{
The choice of norm is not cosmetic. In high dimensions the norms diverge sharply:
for a vector of $10{,}000$ entries each equal to $1$, the $2$-norm is $100$, the
$1$-norm is $10{,}000$, and the $\infty$-norm is $1$. A constraint that is mild in
one may be severe in another, so ``the gradient is bounded'' means nothing until
the norm is named.
}

\whyhere{
This is the last piece the gradient penalty needed. The WGAN-GP objective adds
\[
  \lambda \,\E\Bigl[\bigl(\norm{\nabla D(\hat{x})}_2 - 1\bigr)^2\Bigr]
\]
to the critic's loss, with $\lambda = 10$ throughout this project. Read plainly:
compute the critic's gradient at a point, measure its Euclidean length, and
penalise the square of its deviation from 1. Chapter~\ref{ch:gans} explains why 1
is the target --- it enforces the Lipschitz condition that makes the critic's
output a valid Wasserstein estimate. The gradient-norm clipping applied to every
model here, \texttt{max\_norm=1.0}, is a separate use of the same $2$-norm:
rescale the whole gradient whenever its Euclidean length exceeds 1.
}

%% ─────────────────────────────────────────────────────────────────────────────
\section{Eigenvalues and Eigenvectors}
\label{sec:eigen}

\situation{
Stretch a rubber sheet: twice as wide, three times as tall. Most drawn arrows come
out pointing somewhere new. But an arrow lying exactly horizontal stays
horizontal, merely twice as long; a vertical one stays vertical, three times as
long. Those two directions are special to this particular stretch.
}

\problem{
A general matrix mixes every direction into every other, which makes repeated
application hard to reason about. If we can find the directions the matrix merely
scales, its behaviour becomes transparent: along those directions it is nothing
but multiplication by a number.
}

\workedarith{
Take $A = \begin{bmatrix} 2 & 0 \\ 0 & 3 \end{bmatrix}$.
\[
  A\begin{bmatrix} 1 \\ 0 \end{bmatrix} = \begin{bmatrix} 2 \\ 0 \end{bmatrix}
  = 2\begin{bmatrix} 1 \\ 0 \end{bmatrix} ,
  \qquad
  A\begin{bmatrix} 0 \\ 1 \end{bmatrix} = \begin{bmatrix} 0 \\ 3 \end{bmatrix}
  = 3\begin{bmatrix} 0 \\ 1 \end{bmatrix} .
\]
Directions unchanged, lengths scaled by 2 and 3. Eigenvalues 2 and 3.

\medskip
A generic direction is not preserved:
\[
  A\begin{bmatrix} 1 \\ 1 \end{bmatrix} = \begin{bmatrix} 2 \\ 3 \end{bmatrix} ,
\]
which is not a multiple of $[1,1]$ --- the arrow has rotated.

\medskip
\textbf{A less obvious case.} $B = \begin{bmatrix} 3 & 1 \\ 0 & 2 \end{bmatrix}$.

Try $[1,0]$: $B[1,0]^{\!\top} = [3,0]^{\!\top} = 3[1,0]^{\!\top}$. Eigenvalue 3. \checkmark

Try $[1,-1]$: $B[1,-1]^{\!\top} = [3 - 1,\; -2]^{\!\top} = [2,-2]^{\!\top} = 2[1,-1]^{\!\top}$.
Eigenvalue 2. \checkmark

\medskip
Note in both examples the eigenvalues are the diagonal entries --- true for
triangular matrices, not in general.
}

\formaldef{
A non-zero vector $\mathbf{v}$ is an \textbf{eigenvector} of a square matrix $A$
with \textbf{eigenvalue} $\lambda$ if
\[
  A\mathbf{v} = \lambda \mathbf{v} .
\]
The direction is preserved; only the length changes, by the factor $\lambda$
(reversing if $\lambda < 0$).

Eigenvalues solve $\det(A - \lambda I) = 0$. A real matrix need not have real
eigenvalues --- a rotation has none, since it preserves no direction --- but a
\emph{symmetric} real matrix always does.
}

\analogybreak{
Eigenvectors give a complete picture only when there are enough of them to span
the space, which is not guaranteed. Some matrices are \emph{defective} and cannot
be diagonalised at all. The convenient picture --- ``every matrix is just scaling
along some special axes'' --- is a property of a well-behaved subclass, not a
universal truth.
}

\whyhere{
The Hessian of \S\ref{sec:hessian} is symmetric, so its eigenvalues are real, and
they make ``curves upward in every direction'' precise: all eigenvalues positive
means a minimum, all negative a maximum, mixed signs a saddle. Their ratio is the
condition number, which measures how elongated the valley is and therefore how
badly a single global learning rate will perform --- the quantitative version of
the gorge-versus-bowl problem that motivates Adam.
}

%% ─────────────────────────────────────────────────────────────────────────────
\section{Vector Spaces and Linear Transformations}
\label{sec:vectorspaces}

\situation{
Recipes. Combine half a batch of one and two batches of another and you get a
sensible recipe. Combine sound waves and you get a sound; combine forces and you
get a force. In each case adding two objects and scaling one produce objects of
the same kind. The collection is closed under those two operations.
}

\problem{
The same theorems keep reappearing for lists of numbers, for functions, for random
variables. Rather than reproving each time, we identify what these settings share
--- and then prove things once, for anything that shares it.
}

\workedarith{
\textbf{Is the set of vectors $[x, y]$ with $y = 2x$ a vector space?}

Take $[1,2]$ and $[3,6]$, both on the line. Their sum is $[4,8]$, and
$8 = 2 \times 4$. \checkmark Still on the line.

Scale $[1,2]$ by $-5$: $[-5,-10]$, and $-10 = 2 \times (-5)$. \checkmark

Zero: $[0,0]$ satisfies $0 = 2 \times 0$. \checkmark

Closed under both operations and containing zero --- yes.

\medskip
\textbf{Is the set with $y = 2x + 1$?}

Take $[0,1]$ and $[1,3]$, both on the line. Their sum is $[1,4]$, but
$2 \times 1 + 1 = 3 \neq 4$. $\times$ Not closed. Nor does it contain $[0,0]$,
since $2 \times 0 + 1 = 1 \neq 0$.

The offset $+1$ is fatal, and this is exactly why $f(x) = mx + c$ with $c \neq 0$
is called \emph{affine} rather than linear (\S\ref{sec:linear}).

\medskip
\textbf{Is a matrix map linear?} For $A = \begin{bmatrix} 2 & 0 \\ 1 & 3\end{bmatrix}$,
$\mathbf{u} = [1,0]$, $\mathbf{v} = [0,1]$:
\[
  A(\mathbf{u} + \mathbf{v}) = A[1,1]^{\!\top} = [2,4]^{\!\top} ,
\]
\[
  A\mathbf{u} + A\mathbf{v} = [2,1]^{\!\top} + [0,3]^{\!\top} = [2,4]^{\!\top} . \quad\checkmark
\]
}

\formaldef{
A \textbf{vector space} is a set closed under addition and scalar multiplication,
containing a zero element, with each element having an additive inverse, and
satisfying the usual associative, commutative and distributive laws.

A map $T$ between vector spaces is a \textbf{linear transformation} if
\[
  T(\mathbf{u} + \mathbf{v}) = T(\mathbf{u}) + T(\mathbf{v})
  \quad\text{and}\quad
  T(c\mathbf{v}) = c\,T(\mathbf{v}) .
\]
Every linear transformation between finite-dimensional spaces is a matrix
multiplication, once bases are fixed --- which is why matrices and linear maps are
so often treated as the same thing.

A \textbf{basis} is a minimal set of vectors from which every element can be built
by addition and scaling; the number of them is the \textbf{dimension}.
}

\analogybreak{
Almost nothing interesting in this thesis is linear. Activation functions are
deliberately non-linear, GAN losses are non-convex, and the map from noise to a
synthetic return series is wildly non-linear by design. Linear algebra is not a
description of these systems --- it is the local language in which each small piece
is written, and the thing calculus reduces them to one step at a time.
}

\whyhere{
Every layer is a linear map followed by a non-linearity, and this alternation is
what gives a deep network its expressive power: without the non-linearity the
whole stack would collapse to a single matrix (\S\ref{sec:matmul}). The set of
returns a generator can produce is not a vector space --- it is a curved subset of
$\R^{128}$ --- and Chapter~12 (Generative Models) takes up what that means for the
manifold a generative model is trying to learn.
}
